\documentclass[tikz,border=8pt]{standalone}
\usepackage{ctex}
\usepackage{amsmath}
\usetikzlibrary{calc, arrows.meta}

\begin{document}
\begin{tikzpicture}[scale=1.0, thick]

% 坐标轴
\draw[->, thin, gray] (-0.5,0) -- (5.0,0) node[right, font=\small, gray] {$x$};
\draw[->, thin, gray] (0,-3.2) -- (0,3.2) node[above, font=\small, gray] {$y$};
\node[font=\footnotesize, gray] at (-0.25,-0.2) {$O$};

% 抛物线 y^2 = 4x（p=2，焦点 F(1,0)）
\draw[black, thick, domain=-3.0:3.0, samples=100, variable=\t]
  plot ({\t*\t/4}, {\t});

% 焦点 F(1,0)
\coordinate (F) at (1,0);
\filldraw[black] (F) circle (2pt);
\node[font=\small, below right] at (F) {$F(1,0)$};

% 准线 x = -1（辅助参考）
\draw[gray, dashed, thin] (-1,-3.0) -- (-1,3.0);
\node[font=\footnotesize, gray, above] at (-1,3.0) {准线};

% 4 条平行入射光（水平，从左向右）+ 对应反射光
% y值：-2.4, -1.2, 1.2, 2.4
% 对应抛物线上点：(y^2/4, y)
% y=-2.4: x=1.44  → P1=(1.44,-2.4)
% y=-1.2: x=0.36  → P2=(0.36,-1.2)
% y= 1.2: x=0.36  → P3=(0.36, 1.2)
% y= 2.4: x=1.44  → P4=(1.44, 2.4)

\foreach \yval/\xval in {-2.4/1.44, -1.2/0.36, 1.2/0.36, 2.4/1.44}{
  % 入射光（从 x=-1 水平射到抛物线点）
  \draw[red, thick, -{Stealth[length=5pt,width=3pt]}]
    (-1,\yval) -- (\xval,\yval);
  % 点在抛物线上
  \filldraw[black] (\xval,\yval) circle (1.5pt);
  % 反射光（从抛物线点射向焦点）
  \draw[red, dashed, thick, -{Stealth[length=5pt,width=3pt]}]
    (\xval,\yval) -- (F);
}

% 标注
\node[red, font=\footnotesize] at (-0.5, 3.0) {平行入射};
\draw[red, thin, -{Stealth[length=4pt,width=2.5pt]}] (-1,2.8) -- (-0.1,2.8);

\node[font=\footnotesize] at (3.0, 2.2) {$y^2=4x$};

% 反射汇聚说明
\node[font=\footnotesize, align=center] at (3.5, -2.5)
  {平行光经抛物面\\反射后汇聚于焦点};

\end{tikzpicture}
\end{document}
